%HOMEWORK template for M 325K
\documentclass{article}
\headheight=7pt         \topmargin=14pt
\textheight=574pt       \textwidth=445pt
\oddsidemargin=18pt     \evensidemargin=18pt

%Begin package import

\usepackage{amsmath, amssymb, amsthm}
\usepackage{mathtools}
\usepackage{pdfpages}
\usepackage{tikz-cd}

%End package import
%Begin custom commands

\newcommand{\C}{\mathbb{C}}
\newcommand{\R}{\mathbb{R}}
\newcommand{\Q}{\mathbb{Q}}
\newcommand{\Z}{\mathbb{Z}}
\newcommand{\N}{\mathbb{N}}
\newcommand{\abs}[1]{\left\vert #1\right \vert}
\newcommand{\RM}[1]{\mathrm{#1}}
\newcommand{\BF}[1]{\textbf{#1}}
\newcommand{\e}{\epsilon}
\newcommand{\ceil}[1]{\lceil{#1}\rceil}
\newcommand{\floor}[1]{\lfloor{#1}\rfloor}
\newcommand{\rarr}[0]{\rightarrow}
\newcommand{\IM}[1]{\text{Im}(#1)}
\newcommand{\Hom}[0]{\text{Hom}}
\newcommand{\Pow}[1]{\text{Pow}(#1)}
\newcommand{\angles}[1]{\langle #1 \rangle}

%End custom commands
%Begin custom theorems

\newtheorem{innercustomgeneric}{\customgenericname}
\providecommand{\customgenericname}{}
\newcommand{\newcustomtheorem}[2]{%
\newenvironment{#1}[1]
{%
\renewcommand\customgenericname{#2}%
\renewcommand\theinnercustomgeneric{##1}%
\innercustomgeneric%
}
{\endinnercustomgeneric}
}

\newcustomtheorem{THRM}{Theorem}
\newcustomtheorem{LEM}{Lemma}
\newcustomtheorem{EX}{Exercise}
\newcustomtheorem{COR}{Corollary}
\newcustomtheorem{PROB}{Problem}
\newcustomtheorem{claim}{Claim}

\newenvironment{solution}
{\renewcommand\qedsymbol{$\blacksquare$}\begin{proof}[Solution]}
{\end{proof}}

\newenvironment{definition}
{\renewcommand\qedsymbol{$\blacksquare$}\begin{proof}[Definition]}
{\end{proof}}

%End custom theorems
\begin{document}

\title{Homework 6}
\author{Arham Lodha}%put your name here
\date{February 27, 2024}%enter the due date here
\maketitle

\begin{PROB}{1a}\label{Problem 1a.}
\end{PROB}
\begin{proof}
    For $\epsilon > 0$, by the Archemedian Property, $\exists K \in \N$ such that $K > \frac{\sqrt{5}}{\sqrt{4\epsilon}}$ thus $\epsilon > \frac{5}{4K^2}$ . Then for any $n \geq K$:
    
    \begin{align*}
        \abs{\frac{n^2 - 1}{2n^2 + 3} - \frac{1}{2}} = \abs{\frac{n^2 - 1 - \frac{1}{2}(2n^2 + 3)}{2n^2 + 3}} = \abs{\frac{\frac{-5}{2}}{2n^2 + 3}} &= \abs{-\frac{5}{4n^2 + 6}} \\
        &= \frac{5}{4n^2 + 6} \text{ since $n > 0 \implies 4n^2 + 6 > 0$ and $|-5| = 5$} \\
        &\leq \frac{5}{4n^2} \\
        &\leq \frac{5}{4K^2} \text{ since $n \geq K$ } \\
        &< \epsilon \text{ since $\epsilon > \frac{5}{4K^2}$ } 
    \end{align*}

    Thus $\lim \left(\frac{n^2 - 1}{2n^2 + 3}\right) = 0$.  

\end{proof}

\begin{PROB}{5b}\label{Problem 5b.}
\end{PROB}
\begin{proof}
    For $\epsilon > 0$, by the Archemedian Property, $\exists K \in \N$ such that $K > \frac{1}{\epsilon}$ thus $\epsilon > \frac{1}{K}$. Then for any $n \geq K$, 
    
    \begin{align*}
        \abs{\frac{(-1)^n n}{n^2 + 3} - 0} = \abs{\frac{(-1)^n n}{n^2 + 3}} &= \abs{(-1)^n} \abs{\frac{n}{n^2 + 3}} \\
        &= \abs{\frac{n}{n^2 + 3}} \\
        &\leq \abs{\frac{n}{n^2}} = \abs{\frac{1}{n}} = \frac{1}{n} \\
        &\leq \frac{1}{K} \text{ since $n \geq K$} \\
        &< \epsilon \text{ since $\epsilon > \frac{1}{K}$ }
    \end{align*}

    Thus $\lim \left(\frac{(-1)^n n}{n^2 + 3}\right) = 0$.  

\end{proof}

\bigskip

\begin{PROB}{2}\label{Problem 2.}
    
\end{PROB}
\begin{proof}
    Let $(x_n)$ be a sequence where $lim x_n = x > 0$.   
    Let $\epsilon = \frac{x}{20}$. Thus since $x > 0$, $\epsilon < \frac{x}{2} < x$ and $-\frac{x}{2} < -\epsilon$ and $\epsilon < x$. By definition of limit, for this $\epsilon$, $\exists K \in \N \ni \forall n \geq K$, $-\epsilon < x_n - x < \epsilon$. Thus since $-\frac{x}{2} < -\epsilon$ and $\epsilon < x$, $-\frac{x}{2} < x_n - x < x$. Thus $-\frac{x}{2} + x = \frac{x}{2} < x_n < x + x =2x$, thus $\frac{x}{2} < x_n < 2x$.            
\end{proof}

\bigskip

\begin{PROB}{3a}\label{Problem 3a.}
\end{PROB}
\begin{proof}
    Base case: $n = 1$, $1 < 2$. 
    Induction Hypothesis: $\forall n \in \N$, $n < 2^n$
    Inductive Step: By inductive hypothesis, $2n < 2^{n + 1}$. Since $n \in \N$, $n \geq 1$ thus $n + 1 \leq n + n = 2n$. Thus $n + 1 \leq 2^{n + 1}$. 
    
    Thus by induction, $\forall n \in \N$, $2^n > n$.    

    $\forall h \in \R \ni h > 0$, let $m = \ceil h$, thus $m \geq h$ and $m \in \N$. Thus $m < 2^m$, thus $h < 2^m$. Thus $h$ is not a upperbound. Thus $x_n = 2^n$ has no upper bound and thus is unbounded. By the contrapositive of the theorem that convergent sequences are bounded, $x_n$ is unbounded thus divergent, not convergent.         
\end{proof}

\begin{PROB}{3b}\label{Problem 3b.}
\end{PROB}
\begin{proof}
    Let $a_n = \frac{1}{n}$, $b_n = \frac{1}{n^2}$, $c_n = 4$, $d_n = 2$ .
    $\forall \epsilon > 0$, by Archemedian property, $\exists K, L \in \N \ni K > \frac{1}{\epsilon}$ and $L > \frac{1}{\sqrt{\epsilon}}$, thus $\frac{1}{L^2} < \epsilon > \frac{1}{K}$. Also let $M = 1$.    Thus $\forall k \geq K, l \geq L, m \geq M$, 
    
    \begin{align*}
        \abs{a_n - 0} = \abs{\frac{1}{k} - 0} &= \frac{1}{k} \\
        &\leq \frac{1}{K} \text{ since $k \leq K$} \\
        &< \epsilon \text{ since $\frac{1}{K} < \epsilon$ } 
    \end{align*}

    \begin{align*}
        \abs{b_n - 0} = \abs{\frac{1}{l^2} - 0} &= \frac{1}{l^2} \\
        &\leq \frac{1}{L^2} \text{ since $l \leq L$} \\
        &< \epsilon \text{ since $\frac{1}{L^2} < \epsilon$ } 
    \end{align*}

    \begin{align*}
        \abs{c_n  - 4} = \abs{4 - 4} = 0 < \epsilon 
    \end{align*}

    Thus $\lim a_n = 0$, $\lim b_n = 0$, $\lim c_n = 4$.
    
    Let $x_n = (2 + \frac{1}{n})^2 = 4 + \frac{4}{n} + \frac{1}{n^2} = c_n + 4a_n + b_n$
    
    Since scalar multiple of a sequence which converges also converges and the limit of a scalar multiple of a sequence is equal to the scalar multiple of the limit of the sequence, $\lim 4a_n = 4 \lim a_n = 0$. Since the sum of sequences which converge and their limit is equal to the sum of the limits, $x_n$ conveges because $c_n$, $4a_n$, and $b_n$ converge, and $\lim x_n = \lim c_n + \lim 4a_n + \lim b_n  = 4 + 0 + 0 = 4$.       
\end{proof}

\begin{PROB}{3c}\label{Problem 3c.}
\end{PROB}
\begin{proof}
    $\forall \epsilon > 0$, by the Archemedian property, $\exists K \in \N \ni K > \frac{2}{\epsilon^2}$ thus $\epsilon > \frac{2}{\sqrt{K}}$. Thus $\forall n \geq K$, 
    
    \begin{align*}
        \abs{x_n - 1} = \abs{\frac{\sqrt{n} - 1}{\sqrt{n} + 1} - 1} = \abs{\frac{\sqrt{n} - 1 - \sqrt{n} - 1}{\sqrt{n} + 1}} &= \abs{-\frac{2}{\sqrt{n} + 1}} \\
        &= \frac{2}{\sqrt{n} + 1} \\
        &< \frac{2}{\sqrt{n}} \text{ since $\sqrt{n} < \sqrt{n} + 1$} \\
        &\leq \frac{2}{\sqrt{K}} \text{ since $K \leq n$} \\
        &< \epsilon \text{ since $\frac{1}{\sqrt{K}} < \epsilon$ }
    \end{align*}

    Thus $\lim x_n = 1$ 
\end{proof}

\begin{PROB}{3d}\label{Problem 3d.}
\end{PROB}
\begin{proof}
    Let $x_n = (-1)^n n^2$ 
    $\forall h > 0$, let $m = \ceil h$. If $m \mod 2 = 0$, $m > 1$ thus $m^2 > m$ and $(-1)^m m^2 = m^2 > m \geq h$. Thus $h$ is not a upper bound. Otherwise let $n := m + 1$, thus $n \mod 2  = 0$ and $n >1$, thus $n^2 > n$ and $(-1)^n n^2 = n^2 > n > m \geq h$ Thus $h$ is not a upper bound. Thus $x_n$ is not bounded above and thus is unbounded. Thus by the contrapositive of the theorem that convergent sequences are bounded, since $x_n$ is unbounded it is not convegent and thus is divergent.                  
\end{proof}

\begin{PROB}{3e}\label{Problem 3e.}
\end{PROB}
\begin{proof}
    Let $x_n = \frac{(-1)^n}{n + 2}$. $\forall \epsilon > 0$, by the Archemedian property, $\exists K \in \N \ni K > \frac{1}{\epsilon}$, thus $\epsilon > \frac{1}{K}$. Thus $\forall n \geq K$, 
    
    \begin{align*}
        \abs{x_n - 0} = \abs{\frac{(-1)^n}{n + 2}} &= \frac{1}{n + 2} \\
        &< \frac{1}{n} \text{ since $n < n + 2$} \\
        &\leq \frac{1}{K} \text{ since $n \geq K$ } \\
        &< \epsilon \text{ since $1/K < \epsilon$ }
    \end{align*}

    Thus $\lim x_n = 0$ 
\end{proof}

\end{document}